Put \(s=c^\top\Sigma c>0\) and \(Q=c^\top Vc\).  Theorem \(\mathrm{thm:basin\text{-}stable\text{-}ipwz}\), the assumed weak convergence of the covariance laws, and continuity on the common compact positive-definite set imply
\[
 c^\top\sqrt{T_\ell}(\widehat\theta_{T_\ell}-\theta^*)
 \Rightarrow \sqrt Q\,Z_1,                                                \tag{31}
\]
where \(Z_1\) is independent of \(Q\).  The boundary of every centered finite interval has zero probability under this mixture.  Therefore the definition of \(C_{T_\ell}^{\Sigma}(c,\alpha)\) and (31) give, with \(z=z_{1-\alpha/2}>0\),
\[
 \lim_{\ell\to\infty}P_\ell\{\tau\in C_{T_\ell}^{\Sigma}(c,\alpha)\}
 =P\{|\sqrt QZ_1|\le z\sqrt s\}
 =E\!\left[2\Phi\!\left(z\sqrt{s/Q}\right)-1\right].                 \tag{32}
\]
If \(Q=s\) almost surely, (32) equals \(2\Phi(z)-1=1-\alpha\) for every \(\alpha\), proving sufficiency.

Conversely, suppose the left side of (32) equals \(1-\alpha\) for every \(\alpha\in(0,1)\).  As \(\alpha\) ranges over \((0,1)\), \(z_{1-\alpha/2}\) ranges over \((0,\infty)\).  Hence \(|\sqrt QZ_1|\) and \(|\sqrt sZ_1|\) have the same distribution.  Both underlying variables are symmetric, so \(\sqrt QZ_1\) and \(\sqrt sZ_1\) also have the same distribution.  The compact covariance support bounds \(Q\), and therefore all fourth moments exist.  Equality of the second and fourth moments gives
\[
 E(Q)=s,
 \qquad
 3E(Q^2)=3s^2.
\]
Consequently \(E\{(Q-s)^2\}=0\), or \(c^\top Vc=c^\top\Sigma c\) \(\mu\)-almost surely.  This proves necessity.  The use of every \(\alpha\), equivalently every positive \(z\), is exactly what identifies the absolute-value distribution; no conclusion here is drawn from a single nominal level.